\documentclass[11pt,aspectratio=169]{beamer}

\usetheme{Madrid}
\usecolortheme{seahorse}

\usepackage[utf8]{inputenc}
\usepackage{amsmath,amssymb}
\usepackage{algpseudocode}
\usepackage{algorithmicx}
\usepackage{xcolor}
\usepackage{listings}
\usepackage{booktabs}
\usepackage{tikz}
\usetikzlibrary{arrows.meta,positioning,fit,backgrounds}

\setbeamertemplate{navigation symbols}{}
\setbeamertemplate{footline}[frame number]

\lstdefinestyle{py}{
  language=Python,
  basicstyle=\ttfamily\scriptsize,
  keywordstyle=\color{blue!70!black}\bfseries,
  commentstyle=\color{gray},
  stringstyle=\color{green!40!black},
  showstringspaces=false,
  breaklines=true,
  numbers=left,
  numberstyle=\tiny\color{gray},
  frame=single,
  xleftmargin=1.5em,
  aboveskip=0.3em,
  belowskip=0.3em
}

\title{LeetCode Deep Dive}
\subtitle{Expected Time vs.\ Worst Case: Randomized Selection and Hashing in Practice\\
\small Week 15 Supplement -- TOAE209/TOAH209: Design and Analysis of Algorithms}
\author{Foreign Trade University}
\date{}

\begin{document}

% ---------------------------------------------------------------
\begin{frame}
  \titlepage
\end{frame}
% ---------------------------------------------------------------

\begin{frame}{Outline}
  \tableofcontents
\end{frame}

% =================================================================
\section{A Framework for Attacking a New Problem}
% =================================================================

\begin{frame}{Why walk through LeetCode problems?}
  \begin{itemize}
    \item This week's \texttt{LEETCODE.md} list has eight Medium problems and
    \textbf{no} official Hard -- the same situation as Week 4. This deck picks the
    best-fitting Medium and adds one well-known, correctly-rated Hard as a stretch
    problem, exactly as Week 3 did with \emph{Critical Connections in a Network}.
    \item We pick one \textcolor{blue!50!black}{\textbf{Medium}} (\emph{Kth Largest
    Element in an Array}, via randomized Quickselect) and one
    \textcolor{red!60!black}{\textbf{Hard}} (\emph{Longest Duplicate Substring}, via
    binary search $+$ Rabin--Karp rolling hash) -- both built directly from this week's
    lecture: \textbf{randomized partitioning} and \textbf{universal hashing}.
    \item The crucial new distinction this week teaches, and that both problems will
    force you to be precise about: \textbf{expected} time vs.\ \textbf{worst-case} time
    are \emph{not} the same guarantee.
  \end{itemize}
\end{frame}

\begin{frame}{The four-step process we will repeat twice}
  \begin{enumerate}
    \item \textbf{First instinct.} Write down the algorithm that follows most directly
    from re-reading the problem statement -- usually ``sort everything'' or
    ``compare every pair.''
    \item \textbf{Measure why it hurts.} Compute the actual complexity, plug in the
    problem's constraints, and see the number blow up -- or see that it does
    \emph{more work than the question actually asked for}.
    \item \textbf{Find the structural insight.} A randomized choice (a random pivot, a
    random-looking hash) that lets you avoid the wasted work \emph{in expectation},
    even though no single run is guaranteed to avoid it.
    \item \textbf{Implement the optimal algorithm}, trace it by hand on a small example,
    then look for the \emph{programming-level} tricks that make the implementation
    itself correct and fast.
  \end{enumerate}
  \begin{alertblock}{This mirrors the course}
    Steps 1--3 are exactly this week's ``Las Vegas gambles with time, Monte Carlo
    gambles with correctness'' framing; step 4 is what turns a correct algorithm into
    working code.
  \end{alertblock}
\end{frame}

% =================================================================
\section{Problem 1 (Medium): Kth Largest Element in an Array}
% =================================================================

\begin{frame}{Problem statement}
  \begin{block}{LeetCode 215 -- Kth Largest Element in an Array
  (\textcolor{blue!50!black}{Medium})}
    Given an integer array \texttt{nums} and an integer $k$, return the $k$-th
    \textbf{largest} element (not the $k$-th distinct element).
  \end{block}
  \begin{itemize}
    \item Constraints (LeetCode): $1 \le k \le \texttt{nums.length} \le 10^5$,
    $-10^4 \le \texttt{nums[i]} \le 10^4$.
    \item The $k$-th largest of $n$ elements is exactly the $(n-k)$-th smallest in
    $0$-indexed sorted order -- \textbf{this week's Quickselect}, computing one order
    statistic instead of a full sort.
  \end{itemize}
\end{frame}

\begin{frame}{Worked example}
  \small
  \texttt{nums = [3,2,1,5,6,4]},\quad $k=2$ (2nd largest).
  \begin{itemize}
    \item Sorted ascending: $[1,2,3,4,5,6]$; sorted descending: $[6,5,4,3,2,1]$.
    \item The 2nd largest is $5$ -- equivalently, the element at $0$-indexed position
    $n-k = 6-2 = 4$ in ascending sorted order (value $5$, at index $4$).
    \item \textbf{Answer:} $5$.
  \end{itemize}
  \begin{center}
  \begin{tabular}{@{}c|cccccc@{}}
    index (asc.\ sorted) & 0 & 1 & 2 & 3 & \textbf{4} & 5 \\
    \midrule
    value & 1 & 2 & 3 & 4 & \textbf{5} & 6 \\
  \end{tabular}
  \end{center}
\end{frame}

\begin{frame}{First instinct: sort everything}
  \begin{block}{Most direct reading of the problem}
    ``Return the $k$-th largest'' $\Rightarrow$ sort \texttt{nums}, then read off
    \texttt{nums[n-k]} (ascending) or \texttt{nums[k-1]} (descending).
  \end{block}
  \begin{algorithmic}[1]
    \Function{KthLargestNaive}{nums, $k$}
      \State sort \texttt{nums} in ascending order
      \State \Return $\texttt{nums}[\,n-k\,]$
    \EndFunction
  \end{algorithmic}
  Cost: $O(n\log n)$ -- correct, and often exactly what a candidate writes first.
\end{frame}

\begin{frame}{Why this is wasteful (not wrong, wasteful)}
  \begin{itemize}
    \item Sorting computes \emph{every} order statistic -- the 1st largest, 2nd
    largest, \dots, $n$-th largest -- when the problem only asked for \emph{one} of
    them.
    \item A min-heap of size $k$ improves this to $O(n\log k)$ (push each element,
    pop when size $>k$, answer is the heap's minimum) -- better, but still pays a
    $\log$ factor for information you don't need.
    \item Both approaches impose a total order on data where the question only needs
    ``one specific rank.''
  \end{itemize}
  \begin{alertblock}{The smell to notice}
    Whenever a problem asks for \emph{one} order statistic (a median, a $k$-th
    smallest/largest), full sorting is doing $\log n$ times more work than necessary.
    That gap is exactly what randomized Quickselect closes.
  \end{alertblock}
\end{frame}

\begin{frame}{The structural insight}
  \begin{itemize}
    \item Partitioning around a pivot (this week's lecture) tells you, in one pass,
    \emph{which side} the answer is on -- without needing to know the relative order of
    everything on that side.
    \item So recurse into \textbf{only one} side (unlike QuickSort, which recurses into
    both), and a \textbf{random} pivot keeps the expected shrink factor bounded on
    \emph{every} input, with no adversarial worst case for a fixed pivot rule.
    \item This week's Theorem (Quickselect) proved this gives \textbf{expected}
    $O(n)$ -- linear, not $O(n\log n)$, because you stop paying for the side you never
    recurse into.
  \end{itemize}
\end{frame}

\begin{frame}{Optimal algorithm: partition step (Lomuto partition)}
  \footnotesize
  \begin{algorithmic}[1]
    \Function{Partition}{$A$, left, right, pivotIdx}
      \State swap $A[\text{pivotIdx}]$ and $A[\text{right}]$;\ pivot $\gets A[\text{right}]$
      \State store $\gets$ left
      \For{$i \gets$ left \textbf{to} right$-1$}
        \If{$A[i] <$ pivot} swap $A[i], A[\text{store}]$;\ store $\gets$ store$+1$ \EndIf
      \EndFor
      \State swap $A[\text{store}], A[\text{right}]$;\ \Return store
    \EndFunction
  \end{algorithmic}
  Partitions $A[\text{left}..\text{right}]$ around the chosen pivot in place: everything
  $<$ pivot ends up left of \texttt{store}, the pivot lands exactly at \texttt{store}.
\end{frame}

\begin{frame}{Optimal algorithm: the selection driver}
  \footnotesize
  \begin{algorithmic}[1]
    \Function{Select}{$A$, left, right, target}
      \If{left $=$ right} \Return $A[\text{left}]$ \EndIf
      \State pivotIdx $\gets$ random integer in $[\text{left}, \text{right}]$
      \State store $\gets \Call{Partition}{A, \text{left}, \text{right}, \text{pivotIdx}}$
      \If{target $=$ store} \Return $A[\text{store}]$
      \ElsIf{target $<$ store} \Return $\Call{Select}{A,\text{left},\text{store}-1,\text{target}}$
      \Else\ \Return $\Call{Select}{A,\text{store}+1,\text{right},\text{target}}$
      \EndIf
    \EndFunction
  \end{algorithmic}
  Expected $O(n)$ time -- this week's bound, with target $= n-k$. Only \textbf{one}
  recursive call ever happens (unlike QuickSort's two).
\end{frame}

\begin{frame}{Trace on the worked example}
  \scriptsize
  \texttt{nums = [3,2,1,5,6,4]}, target $= n-k = 4$. (One concrete run; any pivot
  sequence gives the same final answer.)
  \begin{center}
  \begin{tabular}{@{}c p{4.2cm} c p{3.5cm}@{}}
    \toprule
    call & array state before & pivot chosen & result \\
    \midrule
    Select(0,5,4) & $[3,2,1,5,6,4]$ & index $5$ (value $4$) & partition $\to$
    $[3,2,1,4,6,5]$, pivot lands at index $3$; target $4>3$ \\
    Select(4,5,4) & $[6,5]$ (indices $4,5$) & index $5$ (value $5$) & partition $\to$
    $[5,6]$ at indices $4,5$; pivot lands at index $4$; target $4=4$: \textbf{return
    $A[4]=5$} \\
    \bottomrule
  \end{tabular}
  \end{center}
  \begin{exampleblock}{Answer}
    Two recursive calls, one side discarded each time -- returns $5$, matching the
    2nd-largest element. Notice the first call never looked at the relative order of
    indices $0,1,2$ once it knew the answer was on the other side.
  \end{exampleblock}
\end{frame}

\begin{frame}{Complexity: naive vs.\ optimal}
  \begin{center}
  \begin{tabular}{@{}lll@{}}
    \toprule
    Approach & Time & Guarantee \\
    \midrule
    Sort everything & $O(n\log n)$ & worst-case \\
    Min-heap of size $k$ & $O(n\log k)$ & worst-case \\
    \textbf{Randomized Quickselect} & $\boldsymbol{O(n)}$ & \textbf{expected} --
    worst-case $O(n^2)$ \\
    Deterministic (median-of-medians) & $O(n)$ & worst-case, but larger constants \\
    \bottomrule
  \end{tabular}
  \end{center}
  \begin{alertblock}{The precise claim, stated carefully}
    Randomized Quickselect is expected-linear on \emph{every} input, with no
    adversarial case for a fixed pivot rule -- but any \emph{single} run can, with tiny
    probability, still take $\Theta(n^2)$. This week's advances deck covers the fully
    deterministic worst-case-$O(n)$ alternative (Blum--Floyd--Pratt--Rivest--Tarjan,
    1973) if you need a hard guarantee instead of an expected one.
  \end{alertblock}
\end{frame}

\begin{frame}[fragile]{Python implementation (part 1: partition)}
  \begin{lstlisting}[style=py]
import random

def findKthLargest(nums, k):
    target = len(nums) - k          # 0-indexed rank of the k-th largest

    def partition(left, right, pivot_index):
        pivot = nums[pivot_index]
        nums[pivot_index], nums[right] = nums[right], nums[pivot_index]
        store = left
        for i in range(left, right):
            if nums[i] < pivot:
                nums[i], nums[store] = nums[store], nums[i]
                store += 1
        nums[right], nums[store] = nums[store], nums[right]
        return store
  \end{lstlisting}
\end{frame}

\begin{frame}[fragile]{Python implementation (part 2: select)}
  \begin{lstlisting}[style=py]
    def select(left, right):
        if left == right:
            return nums[left]
        pivot_index = random.randint(left, right)
        store = partition(left, right, pivot_index)
        if target == store:
            return nums[store]
        if target < store:
            return select(left, store - 1)
        return select(store + 1, right)

    return select(0, len(nums) - 1)
  \end{lstlisting}
  \vspace{0.5em}
  \texttt{select} is nested inside \texttt{findKthLargest} purely so it can close over
  \texttt{nums} and \texttt{target} without passing them explicitly each call.
\end{frame}

\begin{frame}{Implementation tips \& tricks (the programming side)}
  \small
  \begin{itemize}
    \item \textbf{Select is tail-recursive -- turn it into a loop.} Unlike QuickSort,
    Quickselect only ever recurses into \emph{one} side. That means the recursive call
    can be replaced by simply updating \texttt{left}/\texttt{right} in a
    \texttt{while} loop, eliminating recursion depth entirely -- important since an
    adversarial (if unlucky) sequence of pivot draws can still recurse $\Theta(n)$
    deep, risking a \texttt{RecursionError} at $n$ near Python's default limit.
    \item \textbf{Partition in place, not with new lists.} The lecture's proof-friendly
    pseudocode builds new lists $L=\{y<x\}$, $R=\{y>x\}$; a real implementation
    partitions the \emph{same} array in place (as above) to avoid $O(n)$ allocations
    per recursive call.
    \item \textbf{\texttt{random.randint}, not a hand-rolled generator.} Python's
    default \texttt{random} module (Mersenne Twister) is fast and statistically good
    enough here; there is no reason to reach for \texttt{secrets} (cryptographic
    randomness), which is slower and solves a different problem.
  \end{itemize}
\end{frame}

% =================================================================
\section{Problem 2 (Hard): Longest Duplicate Substring}
% =================================================================

\begin{frame}{Problem statement}
  \begin{block}{LeetCode 1044 -- Longest Duplicate Substring
  (\textcolor{red!60!black}{Hard})}
    Given a string $s$, return any substring that occurs at least twice in $s$
    (occurrences may overlap) with the \textbf{longest} possible length. Return
    \texttt{""} if no such substring exists.
  \end{block}
  \begin{itemize}
    \item Constraints (LeetCode): $2 \le |s| \le 3\times10^4$, lowercase English
    letters only.
    \item This is a \textbf{static} question (unlike Problem 1's \emph{online} order
    statistic) about a whole string -- and it is naturally attacked with this week's
    \textbf{hashing} section, not partitioning.
  \end{itemize}
\end{frame}

\begin{frame}{Worked example}
  \small
  \texttt{s = "banana"} ($n=6$).
  \begin{center}
  \begin{tabular}{@{}c|cccccc@{}}
    index & 0 & 1 & 2 & 3 & 4 & 5 \\
    \midrule
    char & b & a & n & a & n & a \\
  \end{tabular}
  \end{center}
  \begin{itemize}
    \item Substring \texttt{"ana"} occurs at index $1$ (\texttt{s[1:4]}) \textbf{and}
    index $3$ (\texttt{s[3:6]}) -- an \emph{overlapping} duplicate.
    \item No length-$4$ substring repeats (\texttt{"bana"}, \texttt{"anan"},
    \texttt{"nana"} are all distinct).
    \item \textbf{Answer:} \texttt{"ana"} (length $3$).
  \end{itemize}
\end{frame}

\begin{frame}{First instinct: compare every pair of substrings}
  \begin{block}{Most direct reading of the problem}
    For each candidate length $L$ from $n-1$ down to $1$, look at every substring of
    length $L$ and check whether it matches any other substring of the same length.
  \end{block}
  \begin{algorithmic}[1]
    \For{$L \gets n-1$ \textbf{downto} $1$}
      \For{each pair of start indices $i < j$ with $j + L \le n$}
        \If{$s[i..i{+}L) = s[j..j{+}L)$} \Return $s[i..i{+}L)$ \EndIf
      \EndFor
    \EndFor
  \end{algorithmic}
\end{frame}

\begin{frame}{Why this is bad}
  \begin{itemize}
    \item For a fixed $L$, comparing all $O(n)$ pairs of length-$L$ substrings costs
    $O(n)$ comparisons $\times$ $O(L)$ per comparison $= O(nL)$.
    \item Summed over all $n$ possible lengths $L$: $O(n^2 L_{\text{avg}}) = O(n^3)$
    in the worst case.
    \item Plug in the constraint $n = 3\times10^4$: $n^3 = 2.7\times10^{13}$ -- utterly
    infeasible, even though the problem ``only'' asks about a string of length
    $30{,}000$.
  \end{itemize}
  \begin{alertblock}{The smell to notice}
    ``Compare every substring against every other substring'' re-derives string
    equality from scratch each time, throwing away the fact that most substrings share
    huge overlapping prefixes with their neighbors. That redundancy is exactly what a
    \textbf{rolling hash} removes.
  \end{alertblock}
\end{frame}

\begin{frame}{The structural insight}
  \begin{itemize}
    \item \textbf{Monotonicity.} If a duplicated substring of length $L$ exists, one of
    length $L-1$ also exists (trim it) -- so ``does a duplicate of length $\ge L$
    exist'' is a monotone yes/no question in $L$, and \textbf{binary search} finds the
    largest feasible $L$ in $O(\log n)$ probes instead of trying every length.
    \item \textbf{Rolling hash (Rabin--Karp), this week's universal-hashing idea.}
    Treat each length-$L$ window as a base-$26$ number mod a large modulus; sliding the
    window by one position updates the hash in $O(1)$ instead of rehashing $O(L)$
    characters -- so one feasibility check, for a fixed $L$, costs $O(n)$ total.
    \item Combined: $O(\log n)$ binary-search steps $\times\ O(n)$ per step
    $= O(n\log n)$ -- versus the naive $O(n^3)$.
  \end{itemize}
\end{frame}

\begin{frame}{Trace on the worked example}
  \scriptsize
  \texttt{s = "banana"}, $n=6$. Binary search $L$ over $[1, n-1] = [1,5]$.
  \begin{center}
  \begin{tabular}{@{}c c p{7.2cm}@{}}
    \toprule
    iteration & mid $L$ & feasibility check \\
    \midrule
    1 & $3$ & substrings \texttt{ban,ana,nan,ana}: index $3$'s \texttt{"ana"} repeats
    index $1$'s $\Rightarrow$ \textbf{feasible}. Record \texttt{"ana"}; search $L>3$. \\
    2 & $4$ & substrings \texttt{bana,anan,nana}: all distinct $\Rightarrow$
    \textbf{infeasible}; search $L<4$. \\
    -- & -- & range exhausted (lo$=4>$hi$=3$): stop \\
    \bottomrule
  \end{tabular}
  \end{center}
  \begin{exampleblock}{Answer}
    Two feasibility checks (not five, one per length) find the answer \texttt{"ana"} --
    matching the expected output.
  \end{exampleblock}
\end{frame}

\begin{frame}{Complexity: naive vs.\ optimal}
  \begin{center}
  \begin{tabular}{@{}lll@{}}
    \toprule
    Approach & Time & At $n=3\times10^4$ \\
    \midrule
    Compare every substring pair, every length & $O(n^3)$ & $\sim 2.7\times10^{13}$
    ops \\
    \textbf{Binary search $+$ rolling hash} & $\boldsymbol{O(n\log n)}$ &
    $\sim 4.5\times10^5$ ops \\
    \bottomrule
  \end{tabular}
  \end{center}
  \begin{alertblock}{The precise claim, stated carefully}
    A hash \emph{match} is a candidate, not a proof -- two different substrings can
    collide. The implementation below always verifies with a direct string comparison
    on a hash hit, so the algorithm is correct with certainty (Las Vegas in spirit),
    not merely ``probably correct'' (Monte Carlo).
  \end{alertblock}
\end{frame}

\begin{frame}[fragile]{Python implementation (part 1: the feasibility check)}
  \begin{lstlisting}[style=py,aboveskip=0pt,belowskip=0pt,basicstyle=\ttfamily\tiny]
def longestDupSubstring(s):
    n = len(s)
    nums = [ord(c) - ord('a') for c in s]
    base, mod = 26, (1 << 32) - 1

    def search(L):
        h = 0
        for i in range(L):
            h = (h * base + nums[i]) % mod
        seen = {h: [0]}
        power = pow(base, L - 1, mod)          # precomputed once per call
        for start in range(1, n - L + 1):
            h = ((h - nums[start - 1] * power) * base
                 + nums[start + L - 1]) % mod
            if h in seen:
                for j in seen[h]:
                    if s[j:j + L] == s[start:start + L]:
                        return start            # verified, not just hashed
                seen[h].append(start)
            else:
                seen[h] = [start]
        return -1
  \end{lstlisting}
\end{frame}

\begin{frame}[fragile]{Python implementation (part 2: the binary search driver)}
  \begin{lstlisting}[style=py]
    lo, hi, result = 1, n - 1, ""
    while lo <= hi:
        mid = (lo + hi) // 2
        start = search(mid)
        if start != -1:
            result = s[start:start + mid]
            lo = mid + 1                       # duplicate found: try longer
        else:
            hi = mid - 1                       # none found: try shorter
    return result
  \end{lstlisting}
  \vspace{0.5em}
  \texttt{hi = n - 1}, not \texttt{n}: a duplicate needs two \emph{distinct} start
  positions, so the longest possible duplicate has length $n-1$.
\end{frame}

\begin{frame}{Implementation tips \& tricks (the programming side)}
  \small
  \begin{itemize}
    \item \textbf{Precompute the power once per length, not once per position.}
    \texttt{power = pow(base, L-1, mod)} costs $O(\log L)$ via fast modular
    exponentiation and is reused across all $n-L$ rolling updates -- computing it fresh
    inside the loop would silently reintroduce an $O(n\log L)$ factor.
    \item \textbf{Only slice strings on a hash \emph{hit}.} \texttt{s[j:j+L]} is
    $O(L)$; calling it unconditionally on every position would erase the entire benefit
    of rolling hashes. It is called only when two hashes already match -- rare, so the
    amortized cost stays low.
    \item \textbf{A fixed modulus like $2^{32}-1$ is convenient but not adversary-proof.}
    Because it is a well-known constant, a maliciously constructed test string could in
    principle be built to collide against it. This week's own \textbf{universal
    hashing} idea (choose the hash function's parameters at random, or use a large
    random prime modulus / double hashing) removes that risk -- exactly the same fix
    that motivated universal hashing in lecture.
  \end{itemize}
\end{frame}

% =================================================================
\section{Summary}
% =================================================================

\begin{frame}{Summary: two problems, one lesson}
  \footnotesize
  \begin{enumerate}
    \item \textbf{Kth Largest Element} needs only \emph{one} order statistic --
    randomized Quickselect answers it in expected $O(n)$ by recursing into only the
    side that matters, instead of sorting (or heap-ordering) everything.
    \item \textbf{Longest Duplicate Substring} needs a \emph{global} yes/no answer
    about a whole string at many candidate lengths -- monotonicity turns that into a
    binary search, and a rolling hash (this week's universal-hashing idea) makes each
    check $O(n)$ instead of $O(n^2)$.
    \item Both times, the naive instinct did \emph{more work than the question asked
    for} (sorting everything; comparing every substring pair); both times, randomness
    -- a random pivot, a random-looking hash -- let the algorithm skip the unnecessary
    work \emph{in expectation}.
    \item Both times, precision about the guarantee mattered: \textbf{expected} $O(n)$
    for Quickselect (worst case $O(n^2)$, unless you pay median-of-medians' constants),
    and a rolling hash that is correct \emph{with certainty} only because every match is
    verified before being trusted.
  \end{enumerate}
  \begin{alertblock}{This closes out the course's LeetCode-walkthrough series}
    \texttt{LEETCODE.md} lists seven more Week 15 problems (\emph{Shuffle an Array},
    \emph{Random Pick with Weight}, \emph{Insert Delete GetRandom O(1)}, \dots) -- and
    fourteen weeks of problems before them. The same four-step process applies to all
    of them: find the naive approach, measure why it hurts, find the structural
    insight, then mind the implementation details.
  \end{alertblock}
\end{frame}

\end{document}
